\section{Čtení vstupu z ovladačů}

Pro zajištění přesného časování řídicích signálů a hardwarové akcelerace komunikace byl využit
subsystém PIO (\textit{Programmable I/O}), který je specifický pro mikrokontroléry RP2040/RP2350.

\subsection{Implementace}

Inicializace PIO automatu probíhá pomocí dedikovaných funkcí:

\begin{itemize}
    \item \texttt{pio\_claim\_free\_sm\_and\_add\_program()} -- zkompilovaný PIO program nahraje do
        instrukční paměti, zkonfiguruje mapování příslušných GPIO pinů a nastaví děličku hodin.
        Druhá funkce nastaví piny.
    \item \texttt{joypad\_read\_program\_init()} -- nastaví piny na potřebné vývody.
\end{itemize}

Komunikace je rozdělena na hardwarovou část (PIO program) a softwarovou obsluhu v jazyce C. PIO
program (výpis \ref{lst:pio_joypad}) je navržen tak, aby po obdržení signálu od procesoru
vygeneroval impulz na pinu \texttt{LATCH} a následně pomocí osmi hodinových cyklů na pinu
\texttt{CLOCK} paralelně přečetl datové bity z obou posuvných registrů.

Taktování PIO automatu je softwarově sníženo na frekvenci přibližně 1~MHz (souhlasí s datasheetem
posuvného registru CD4021), takže spuštění jedné instrukce trvá 1~$\mu$s.

\begin{listing}[htbp]
\small
\begin{minted}{gas}
; PIO program pro taktování čipu CD4021 a čtení jeho sériového výstupu
.program joypad_read

; Set piny ovládají signál LATCH (Strobe)
; Sideset piny ovládají signál CLOCK
; In piny čtou datové výstupy ovladačů
.side_set 1

; Definice zpoždění (v cyklech) pro dodržení časování CD4021 (~15 us)
.define CLK_TIME 15
.define STR_TIME 15

.wrap_target
    pull block         side 0             ; Čekání na příkaz od CPU (přes TX FIFO)
    mov isr, null      side 0

    set pins, 1        side 0             ; LATCH = 1 (Vzorkování tlačítek)
    nop                side 0 [STR_TIME]  ; Prodleva
    set pins, 0        side 0             ; LATCH = 0

    set x, 7           side 0             ; Nastavení počítadla smyčky na 8 iterací
read_bit:
    in pins, 2         side 0 [CLK_TIME]  ; CLK = 0, přečtení 2 pinů současně
    jmp x-- read_bit   side 1             ; CLK = 1, skok na další bit

    push block         side 0             ; Odeslání nasbíraných dat do RX FIFO
.wrap
\end{minted}
\caption{Zdrojový kód PIO programu pro asynchronní čtení ovladačů.}
\label{lst:pio_joypad}
\end{listing}

Před zahájením emulace každého snímku je volána funkce \texttt{update\_joypads()}. Její obsluha
spočívá v odeslání prázdného povelu (nuly) do TX fronty stavového automatu, čímž se spustí smyčka
PIO programu. PIO následně přečte 8 dvojic bitů (jeden bit pro každý ovladač v každém cyklu) a
16bitový výsledek vrátí přes RX frontu. Datové slovo je následně v jazyce C bitovými operacemi
rozděleno do samostatných bytů pro první a druhý ovladač. Táto funkce také zajišťuje synchronizaci s
PIO blokem.
